\documentclass{article}
\usepackage{anyfontsize}
\usepackage[UTF8]{ctex} % 如果需要支持中文
\usepackage{fontspec} 
\setCJKmainfont{SourceHanSansCN-Light}[
    BoldFont=SourceHanSansCN-Medium
]

\usepackage[a4paper, left=0.1in, right=0.1in, top=0.05in, bottom=0.05in]{geometry} % 设置四边页边距为0.1英寸{geometry} % 设置页边距为0.5英寸
\usepackage{enumitem} % 用于自定义列表
\usepackage{amsmath}  % 数学公式支持

\setlength{\parindent}{0pt} % 不缩进
\usepackage{etoolbox} % 用于列表操作
%调整类代码
\usepackage{comment}
\usepackage{tocloft} % 用于定制目录样式
\usepackage{hyperref} %不需要目录盒子注释这
\usepackage{ifthen}
\usepackage{tikz}
\usepackage{titletoc}
\usepackage{graphicx}
\usepackage{amssymb}  % 提供数学符号，包括\therefore
%目录设定
%快捷键 CMD + / 注释
%  \renewcommand{\cftdot}{} % 隐藏目录中的页码
%  \cftpagenumbersoff{section}
%  \cftpagenumbersoff{subsection}
%  \makeatletter
%  \renewcommand{\@pnumwidth}{0em}  % 设置页码宽度为0
%  \renewcommand{\@tocrmarg}{0em}   % 设置右边距为0
%  \makeatother
 

%\setCJKmainfont[
 %   Path = C:/USERS/11252/APPDATA/LOCAL/MICROSOFT/WINDOWS/FONTS/,
 %   UprightFont = {SourceHanSerifCN-Light-5.otf},
 %   BoldFont = {SourceHanSerifCN-Medium-6.otf},
%    Scale=1
%]{SourceHanSerifCN}

%\setmainfont{SourceHanSerif}  % 设置英文字体

% 处理冲突的命令
\let\oldbackepsilon\backepsilon
\let\backepsilon\relax
\let\oldeth\eth
\let\eth\relax
\let\olddigamma\digamma
\let\digamma\relax

\usepackage[a4paper, left=0.1in, right=0.1in, top=0.05in, bottom=0.05in]{geometry} % 设置四边页边距为0.1英寸{geometry} % 设置页边距为0.5英寸
\usepackage{enumitem} % 用于自定义列表
\usepackage{amsmath}  % 数学公式支持
\usepackage{float} % 用于浮动环境

\setlength{\parindent}{0pt} % 不缩进
\usepackage{etoolbox} % 用于列表操作
%调整类代码
\usepackage{comment}
\usepackage{tocloft} % 用于定制目录 样式
\usepackage{hyperref} %不需要目录盒子注释这
\usepackage{ifthen}
\usepackage{tikz}
\usepackage{titletoc}
\usepackage{graphicx}
\usepackage{amssymb}  % 提供数学符号，包括\therefore
%目录设定
%快捷键 CMD + / 注释
%  \renewcommand{\cftdot}{} % 隐藏目录中的页码
%  \cftpagenumbersoff{section}
%  \cftpagenumbersoff{subsection}
%  \makeatletter
%  \renewcommand{\@pnumwidth}{0em}  % 设置页码宽度为0
%  \renewcommand{\@tocrmarg}{0em}   % 设置右边距为0
%  \makeatother
 


\usepackage{environ}

\newif\ifshowcontent  % 定义一个条件判断变量
\showcontenttrue    % 设置为 false，隐藏所有 question 环境的内容

\newcounter{question} % 题目计数器
\setcounter{question}{0}
\NewEnviron{question}{%
    \ifshowcontent
        \par\stepcounter{question}\textbf{\thequestion.} \BODY\par % 如果显示内容，则输出
    \fi
}

\NewEnviron{choices}{%
    \ifshowcontent
        \begin{enumerate}[label=\Alph*., leftmargin=*]
            \BODY
        \end{enumerate}\par
    \fi
}

\NewEnviron{correctanswer}{%
    \ifshowcontent
        \textbf{答案:}\BODY\par
    \fi
}



\newenvironment{explanation}
{\textbf{解析:}\par}
{\par}



% 新增考察方面命令
\newcommand{\checklist}[1]{
    \textbf{考察:} #1 \par % 在题目下方显示考察方面
    \addcontentsline{toc}{subsection}{题号\thequestion 考察: #1} % 将考察内容添加到目录
    \vspace{2em} % 添加1em的垂直空间
}

\graphicspath{{images/}} %统一


\begin{document}
 %\fontsize{22}{31}\selectfont
 \tableofcontents % 添加目录
\newpage
%\pagestyle{empty}




\section{考点1:相关性基础}
\setcounter{question}{0}
\begin{question}
    Suppose an individual buys a correlation swap with a fixed correlation of 0.3 and a notional value of \$2 million for one year. The realized pairwise correlations of the daily log returns at maturity for three assets are $\rho_{1,2}= 0.8$, $\rho_{3,1} = 0.3$, and $\rho_{3,2} = 0.4$. What is the correlation swap buyer's payoff at maturity?
\end{question}

\begin{choices}
    \item \$200,000
    \item \$400,000
    \item \$600,000
    \item \$800,000
\end{choices}

\begin{correctanswer}
    B
\end{correctanswer}

\begin{explanation}
    \textbf{步骤1：确定参数}
    \begin{itemize}
        \item 名义本金：\$2{,}000{,}000
        \item 固定相关：$0.3$
        \item 实现两两相关：$\rho_{1,2}=0.8,\ \rho_{3,1}=0.3,\ \rho_{3,2}=0.4$
        \item 资产数：$n=3$
    \end{itemize}

    \textbf{步骤2：计算平均实现相关}
    \begin{itemize}
        \item 公式：$\displaystyle \bar{\rho}=\frac{2}{n^2-n}\sum_{i<j}\rho_{i,j}$
        \item 代入：$\displaystyle \bar{\rho}=\frac{2}{3^2-3}\,(0.8+0.3+0.4)=\frac{2}{6}\times1.5=0.5$
    \end{itemize}

    \textbf{步骤3：计算到期支付}
    \begin{itemize}
        \item 买方支付函数：$\text{Payoff}=\text{Notional}\times(\bar{\rho}-\rho_{\text{fixed}})$
        \item 代入：\$2{,}000{,}000 $\times (0.5-0.3)=\$400{,}000$
    \end{itemize}
   
   
\end{explanation}
\checklist{
    相关性互换的支付计算
}



\begin{question}
    Within the framework of risk analysis, which of the following choices would be considered most critical when looking at risks within financial institutions?
\end{question}

\begin{choices}
    \item Computing separate capital requirements for a bank's trading and banking books.
    \item Proper analysis of stressed VaR.
    \item Persistent use of backtesting.
    \item Consideration of interactions among risk factors.
\end{choices}

\begin{correctanswer}
    D
\end{correctanswer}

\begin{explanation}

    \textbf{A选项错误。}
    \begin{itemize}
        \item 交易账簿与银行账簿分别计提资本是监管规定与核算边界问题，本身并不解决风险度量中最核心的识别与传导问题。
        \item 即便分账计提充分，仍可能忽视账簿之间与风险类型之间的联动与传染。
    \end{itemize}

    \textbf{B选项错误。}
    \begin{itemize}
        \item 压力VaR是重要环节，但属于“单一风险度量在极端场景下的放大镜”，仍多为单维视角。
        \item 若未建模风险间联动（如市场—信用、市场—流动性），压测结果容易低估“组合放大效应”（如错向风险、流动性螺旋）。
    \end{itemize}

    \textbf{C选项错误。}
    \begin{itemize}
        \item 回测用于检验模型覆盖率与稳定性，侧重事后验证而非事前识别联动机制。
        \item 持续回测不能替代对相关性、尾部依赖与结构性传染渠道的刻画。
    \end{itemize}

    \textbf{D选项正确。}
    \begin{itemize}
        \item 金融机构的关键风险在于风险因子之间的相互作用：相关性时变、尾部依赖（copula/极值相关）、以及机制性传染（如保证金—被动平仓—价差扩大）。
        \item 典型路径：
        \begin{itemize}
            \item 市场风险↑ → 抵押品价值↓ → 融资/流动性风险↑ → 被动减仓 → 市场风险进一步↑（流动性螺旋）
            \item 利率/汇率冲击 → 债务可负担性恶化 → 信用利差走阔（市场—信用联动）
            \item 集中度与错向风险：敞口与违约概率同向（wrong-way risk）放大损失分布尾部
        \end{itemize}
        \item 因此，系统性、传染性与组合效应的识别与量化（情景—相关性—流动性三维一体）是最关键的风险分析任务。
    \end{itemize}

\end{explanation}

\checklist{
    掌握风险分析的关键（因为风险不是独立的，相互作用可能放大或减少总风险）
}


\begin{question}
    The risk measure that extracts the dependence structure from the joint distribution function created from several continuous marginal distribution functions is known as a:
\end{question}

\begin{choices}
    \item Copula
    \item Volatility
    \item Correlation
    \item Multivariate correlation
\end{choices}

\begin{correctanswer}
    A
\end{correctanswer}
\begin{explanation}

    \textbf{A选项正确。}
    \begin{itemize}
        \item Copula基于Sklar定理：将联合分布$F_{X,Y}$分解为边际分布$F_X,F_Y$与依赖结构$C$，即$F_{X,Y}(x,y)=C(F_X(x),F_Y(y))$。
        \item 能在不改变边际的前提下，单独刻画与校准依赖关系（含非线性、非对称、尾部依赖）。
        \item 可灵活选择族（高斯、t、阿基米德族等），适配不同尾部与相关特征。
    \end{itemize}

    \textbf{B选项错误。}
    \begin{itemize}
        \item 波动率度量单变量的离散程度，不描述变量间依赖结构。
    \end{itemize}

    \textbf{C选项错误。}
    \begin{itemize}
        \item 相关性主要反映线性依赖，难以捕捉非线性关系与尾部依赖，信息维度不及Copula。
    \end{itemize}

    \textbf{D选项错误。}
    \begin{itemize}
        \item 多元相关性仍是基于线性相关矩阵的刻画，对非线性与尾部依赖刻画不足，灵活性不及Copula。
    \end{itemize}

\end{explanation}

\checklist{
    掌握Copula的特点（从联合分布中提取依赖结构，可以建模非线性依赖和尾部依赖，比相关性更灵活）
}


\begin{question}
    What type of liquidity risk is most troublesome for complex trading positions?
\end{question}

\begin{choices}
    \item Endogenous
    \item Market-specific
    \item Exogenous
    \item Spectral
\end{choices}

\begin{correctanswer}
    A
\end{correctanswer}

\begin{explanation}

    \textbf{A选项正确。}
    \begin{itemize}
        \item 内生流动性风险源自头寸本身及其交易行为（如规模、集中度、交易频率、参与者结构）。
        \item 对复杂头寸（非线性、分层结构、跨市场联动、定制化合约）更麻烦：一旦减仓会产生较大价格冲击，流动性随卖压恶化，形成自我强化。
        \item 典型表现：冲击成本随成交量上升而非线性放大、订单簿深度在压力下迅速消失、交易对手/做市商收缩报价。
    \end{itemize}

    \textbf{B选项错误。}
    \begin{itemize}
        \item 市场特定（market-specific）流动性风险偏外生，如单一市场/品种在特定时点整体流动性变差，并非复杂头寸独有。
    \end{itemize}

    \textbf{C选项错误。}
    \begin{itemize}
        \item 外生（exogenous）流动性风险来自宏观/监管/系统性冲击，更多是全市场性质，与单个复杂头寸的结构性脆弱性关系较弱。
    \end{itemize}

    \textbf{D选项错误。}
    \begin{itemize}
        \item “谱（spectral）”属于风险度量方法学术语，不是流动性风险分类。
    \end{itemize}

\end{explanation}

\checklist{
    掌握流动性风险类型（内生流动性风险对复杂交易头寸最麻烦，因为它源于头寸本身，随规模增加，难以预测）
}


\begin{question}
    In May of 2008, several large hedge funds had speculative positions in the collateralized debt obligations (CDOs) tranches. These hedge funds were forced into bankruptcy due to the lack of understanding of correlations across tranches. Which of the following statements best describes the positions held by hedge funds at this time and the role of changing correlations? Hedge funds held:
\end{question}

\begin{choices}
    \item a long equity tranche and short mezzanine tranche when the correlations in both tranches decreased.
    \item a short equity tranche and long mezzanine tranche when the correlations in both tranches increased.
    \item a short senior tranche and long mezzanine tranche when the correlation in the mezzanine tranche decreased.
    \item a long mezzanine tranche and short equity tranche when the correlation in the mezzanine tranche increased.
\end{choices}

\begin{correctanswer}
    A
\end{correctanswer}

\begin{explanation}

    \textbf{A选项正确。}
    \begin{itemize}
        \item 头寸结构：做多权益层（Equity Tranche），做空中间层（Mezzanine Tranche）。
        \item 相关性下降的影响（基于高斯Copula分层特性）：
        \begin{itemize}
            \item 权益层：相关性下降→违约更分散、低强度亏损更常见→权益层更容易被持续“磨损”→价值下降（多头亏损）。
            \item 中间层：相关性下降→极端簇拥违约概率下降→中间层受损减少→价值上升（空头亏损）。
        \end{itemize}
        \item 结果：两腿同时亏损（权益多头亏、夹层空头亏），若加杠杆易引发重大损失与被动清算。
    \end{itemize}

    \textbf{B选项错误。}
    \begin{itemize}
        \item 若做空权益、做多夹层且相关性上升，则通常权益受益（价值上升）而夹层受损（价值下降），并非导致对冲基金普遍亏损的组合。
    \end{itemize}

    \textbf{C选项错误。}
    \begin{itemize}
        \item 与题干描述不符：并非做空优先、做多夹层；且仅夹层相关性变化无法解释两腿同向亏损。
    \end{itemize}

    \textbf{D选项错误。}
    \begin{itemize}
        \item 做多夹层、做空权益且夹层相关性上升时：夹层价值下降（多头亏），权益价值通常上升（空头亏），并非题干的“因相关性误判导致双向受损”的经典情形。
    \end{itemize}

\end{explanation}

\checklist{
    掌握CDO分层交易策略（做多权益层和做空中间层的组合，相关性下降会导致双重损失）
}


\begin{question}
    Suppose a creditor makes a \$6 million loan to company A and a \$6 million loan to company B. Based on historical information of companies in this industry, companies A and B each have a 7\% default probability and a default correlation coefficient of 0.6. The expected loss for this creditor under the worst case scenario assuming loss given default is 100\% is closest to:
\end{question}

\begin{choices}
    \item \$420,000
    \item \$504,000
    \item \$588,000
    \item \$840,000
\end{choices}

\begin{correctanswer}
    B
\end{correctanswer}

\begin{explanation}

   
    \textbf{步骤1：已知参数}
    \begin{itemize}
        \item 贷款规模：\$6{,}000{,}000 给A，\$6{,}000{,}000 给B，总额\$12{,}000{,}000
        \item 单体违约概率：$P(A)=P(B)=0.07$
        \item 违约相关系数：$\rho=0.6$
        \item LGD：100\%
    \end{itemize}

    \textbf{步骤2：近似计算联合违约概率（违约指示变量相关近似）}
    \[
\begin{aligned}
P(A\cap B)
&\approx p_A p_B + \rho\,\sqrt{p_A(1-p_A)\,p_B(1-p_B)} \\
&= 0.07\times 0.07 \;+\; 0.6 \times \sqrt{0.07\times 0.93 \times 0.07 \times 0.93} \\
&\approx 0.042 \\[6pt]
\end{aligned}
\]

    \textbf{步骤3：计算“最糟情形”下的预期损失（同时违约的期望损失）}
    \[
\begin{aligned}
\text{EL}_{\text{worst}}
&= P(A\cap B)\times \text{EAD}_{A+B}\times \text{LGD} \\
&\approx 0.042 \times \$12{,}000{,}000 \times 100\% \\
&= \$504{,}000
\end{aligned}
\]
   
\end{explanation}

\checklist{
    掌握联合违约概率计算（考虑违约概率和违约相关性，计算最坏情况下的预期损失）
}


\begin{question}
    New copula correlation models were used by traders and risk managers during the 2008-2010 global financial crisis. This led to miscalculations in the underlying risk for structured products such as collateralized debt obligation (CDO) models. Which of the following statements least likely explains the failure of these new copula correlation models during the financial crisis?
\end{question}

\begin{choices}
    \item The copula correlation models assumed a negative correlation between the equity and senior tranches of CDOs.
    \item Correlations for equity tranches of CDOs increased during the financial crisis.
    \item The correlation copula models were calibrated with data from time periods that had low risk.
    \item Correlations for senior tranches of CDOs decreased during the financial crisis.
\end{choices}

\begin{correctanswer}
    D
\end{correctanswer}

\begin{explanation}

    \textbf{A选项错误（可解释模型失败）。}
    \begin{itemize}
        \item 危机前后，分层之间的相关结构被严重误设（如将权益层与优先层的联动方向/强度设错，或用恒定相关替代时变相关）。
        \item 一旦基本相关结构设定不当（方向、幅度、稳定性），在高相关与尾部依赖上升的环境中会系统性误估分层损失与保护强度，从而造成模型失灵。
    \end{itemize}

    \textbf{B选项错误（可解释模型失败）。}
    \begin{itemize}
        \item 危机中权益层（最先承受违约损失）的相关性显著上升，违约更“同向”与“簇拥”，权益层损失大幅加剧。
        \item 相关性上行叠加尾部依赖，导致基于低相关/恒定相关假设的模型显著低估权益层风险，属于失败的重要原因。
    \end{itemize}

    \textbf{C选项错误（可解释模型失败）。}
    \begin{itemize}
        \item 以低风险时期数据进行校准（低波动、低相关、弱尾部依赖）会在危机时系统性失真。
        \item 由此得到的“静态、偏低”的相关输入无法反映危机期间的时变相关与尾部联动，直接导致CDO分层风险被低估。
    \end{itemize}

    \textbf{D选项正确（最不可能解释失败）。}
    \begin{itemize}
        \item 事实恰相反：危机中优先层相关性并非下降而是上升，保护分层之间的联动增强。
        \item 因此“优先层相关性下降”与实际相背，不能解释模型失败，是最不可能的原因。
    \end{itemize}

\end{explanation}

\checklist{
    掌握Copula模型失败原因（假设负相关、使用低风险数据校准、危机期间相关性增加，但优先层相关性也增加而非减少）
}


\begin{question}
    Expected shortfall is generally suggested as a better alternative than VaR during market turmoil. Which of the following statements regarding VaR and ES is true?
\end{question}

\begin{choices}
    \item ES leads to more required economic capital than VaR does for the same confidence level.
    \item While VaR ensures that the estimate of portfolio risk is less than or equal to the sum of the risks of that portfolio's positions, ES does not.
    \item Both VaR and ES account for the severity of losses beyond the confidence threshold.
    \item ES is relatively easier to backtest than VaR.
\end{choices}

\begin{correctanswer}
    A
\end{correctanswer}

\begin{explanation}

    \textbf{A选项正确。}
    \begin{itemize}
        \item ES是尾部期望损失，考虑超过VaR阈值的所有损失的均值，因此通常大于等于VaR。
        \item 在相同置信水平下（如95\%），ES会要求更多的经济资本，因为它更保守地捕捉了尾部风险。
        \item 例如：如果95\% VaR是\$100，95\% ES会计算所有超过\$100的损失的平均值，结果必然≥\$100。
    \end{itemize}

    \textbf{B选项错误。}
    \begin{itemize}
        \item 事实相反：ES满足次可加性（$\text{ES}(X+Y)\leq\text{ES}(X)+\text{ES}(Y)$），而VaR可能违反次可加性。
        \item 因此ES能更好地处理投资组合风险，体现分散化效果；VaR则可能高估组合风险。
    \end{itemize}

    \textbf{C选项错误。}
    \begin{itemize}
        \item VaR只给出置信水平下的分位数阈值，不提供超过阈值后的损失严重程度信息。
        \item 只有ES考虑尾部损失的严重程度，因为它计算的是超过VaR阈值的平均损失。
    \end{itemize}

    \textbf{D选项错误。}
    \begin{itemize}
        \item ES的回测比VaR更困难，因为需要验证尾部平均损失，而不是简单的例外计数。
        \item ES回测需要更多数据（特别是尾部数据），且验证方法更复杂，统计检验的要求更高。
    \end{itemize}

\end{explanation}

\checklist{
    掌握ES和VaR的区别（ES需要更多资本，具有次可加性，考虑尾部损失严重程度，但回测更困难）
}


\begin{question}
    Which of the following measures is most likely an example of a dynamic financial correlation measure?
\end{question}

\begin{choices}
    \item Pairs trading
    \item Value at risk (VaR)
    \item Binomial default correlation model
    \item Correlation copulas for collateralized debt obligations (CDOs)
\end{choices}

\begin{correctanswer}
    A
\end{correctanswer}

\begin{explanation}

    \textbf{A选项正确。}
    \begin{itemize}
        \item 配对交易（Pairs Trading）是动态相关性度量的典型例子。
        \item 策略特征：
        \begin{itemize}
            \item 实时监控两只相关资产的价格相关性（协整关系）
            \item 当相关性偏离历史均值时，做多被低估资产、做空被高估资产
            \item 根据相关性变化动态调整头寸规模和止损点
        \end{itemize}
        \item 动态性体现在：持续监测相关性指标（如相关系数、协整系数），并根据市场条件变化实时调整交易策略，而非使用固定的相关假设。
    \end{itemize}

    \textbf{B选项错误。}
    \begin{itemize}
        \item VaR通常基于固定时间窗口的历史数据或参数假设，相关结构相对静态。
        \item 虽然存在动态VaR（如EWMA、GARCH等），但VaR本身更多是风险度量工具而非“动态相关性度量”的典型例子。
    \end{itemize}

    \textbf{C选项错误。}
    \begin{itemize}
        \item 二项违约相关模型通常假设固定的违约概率和相关性参数，在给定时期内是静态的。
        \item 不捕捉相关性随时间变化的动态特征。
    \end{itemize}

    \textbf{D选项错误。}
    \begin{itemize}
        \item CDO copula模型在金融危机前常使用固定的相关结构（如恒定高斯copula），无法反映相关性的时变特征。
        \item 这正是2008年金融危机中CDO模型失败的重要原因之一。
    \end{itemize}

\end{explanation}

\checklist{
    掌握动态相关性度量（配对交易是动态相关性的例子，因为它实时监控相关性变化并调整策略）
}


\begin{question}
    A portfolio manager owns a portfolio of options on a non-dividend paying stock XYZ. The portfolio is made up of 18,000 deep in-the-money call options on XYZ and 45,000 deep out-of-the-money call options on XYZ. The portfolio also contains 22,000 forward contracts on XYZ. XYZ is trading at USD 95. If the volatility of XYZ is 28\% per year, which of the following amounts would be closest to the 1-day VaR of the portfolio at the 95 percent confidence level, assuming 252 trading days in a year?
\end{question}

\begin{choices}
    \item USD 1,045
    \item USD 104,500
    \item USD 125,400
    \item USD 145,600
\end{choices}

\begin{correctanswer}
    B
\end{correctanswer}

\begin{explanation}

       \textbf{步骤1：确定各头寸的Delta}
        \begin{itemize}
            \item 深度实值看涨期权（18{,}000张）：$\Delta \approx 1$
            \item 深度虚值看涨期权（45{,}000张）：$\Delta \approx 0$
            \item 远期合约（22{,}000份）：$\Delta = 1$
            \item 组合总$\Delta$：$18{,}000\times1 + 45{,}000\times0 + 22{,}000\times1 = 40{,}000$
        \end{itemize}

         \textbf{步骤2：换算为1日波动率}
        \begin{itemize}
            \item 年化波动率$\sigma=28\%$，交易日$T=252$
            \item 1日波动率：$\sigma_{\text{1d}}=\sigma\sqrt{\tfrac{1}{T}}=0.28\sqrt{\tfrac{1}{252}}$
        \end{itemize}

        \textbf{步骤3：计算1日95\% VaR（Delta-正态近似）}
        \[
        \text{VaR}
        = \alpha \times S \times \Delta \times \sigma \times \sqrt{\tfrac{1}{T}}
        = 1.645 \times 95 \times 40{,}000 \times 0.28 \times \sqrt{\tfrac{1}{252}}
        \approx 104{,}500
        \]

 
   

\end{explanation}

\checklist{
    掌握期权组合VaR的计算
}


\begin{question}
    The relationship of correlation risk to credit risk is an important area of concern for risk managers. Which of the following statements regarding default probabilities and default correlations is incorrect?
\end{question}

\begin{choices}
    \item Creditors benefit by diversifying exposure across industries to lower the default correlations of debtors.
    \item The default term structure increases with time to maturity for most investment grade bonds.
    \item The probability of default is higher in the long-term time horizon for non-investment grade bonds.
    \item For non-investment grade bonds, the probability of default is higher in the immediate time horizon.
\end{choices}

\begin{correctanswer}
    C
\end{correctanswer}

\begin{explanation}

    \textbf{A选项错误（陈述正确）。}
    \begin{itemize}
        \item 跨行业分散投资可降低债务人间的违约相关性（不同行业受不同经济因素影响）。
        \item 这有助于降低组合信用风险，因此该陈述正确。
    \end{itemize}

    \textbf{B选项错误（陈述正确）。}
    \begin{itemize}
        \item 投资级债券的违约期限结构通常随到期时间增加而上升（向上倾斜）。
        \item 更长期限意味着更多不确定性，累积违约概率通常更高，因此该陈述正确。
    \end{itemize}

    \textbf{C选项正确（陈述错误）。}
    \begin{itemize}
        \item 非投资级债券的违约曲线通常呈"驼峰形"或"递减形"。
        \item 特征：
        \begin{itemize}
            \item 即期（短期）违约风险很高
            \item 若公司存活过初始高风险期，违约风险随时间推移而降低
            \item 长期违约概率通常低于短期，而非更高
        \end{itemize}
        \item 因此，陈述"非投资级债券在长期时间范围内违约概率更高"是错误的。
    \end{itemize}

    \textbf{D选项错误（陈述正确）。}
    \begin{itemize}
        \item 非投资级债券在即期时间范围内违约概率确实更高，符合实际观察与信用风险模型。
        \item 因此该陈述正确。
    \end{itemize}

\end{explanation}

\checklist{
    掌握违约概率特征
}


\begin{question}
    The dependence structure between the returns of financial assets plays an important role in risk measurement. For liquid markets, which of the following statements is incorrect?
\end{question}

\begin{choices}
    \item Correlation is valid measure of dependence between random variables for only certain types of return distributions.
    \item Even if the return distributions of two assets have a correlation of zero, the returns of these assets are not necessarily independent.
    \item Copulas make it possible to model marginal distributions and the dependence structure separately.
    \item With short time horizons (3 months or less), correlation estimates are typically very stable.
\end{choices}

\begin{correctanswer}
    D
\end{correctanswer}

\begin{explanation}

    \textbf{A选项错误（陈述正确）。}
    \begin{itemize}
        \item 相关性作为依赖度量主要适用于线性关系和椭圆分布（如正态分布、t分布）。
        \item 对于非线性依赖、非对称分布或存在尾部依赖的情况，相关性可能无法完整刻画依赖结构。
        \item 因此该陈述正确。
    \end{itemize}

    \textbf{B选项错误（陈述正确）。}
    \begin{itemize}
        \item 零相关只意味着不存在线性关系，但不意味着独立。
        \item 可能存在非线性依赖（如二次关系、条件依赖）或尾部依赖，即使相关系数为零。
        \item 因此该陈述正确。
    \end{itemize}

    \textbf{C选项错误（陈述正确）。}
    \begin{itemize}
        \item Copula的核心优势在于能够将边际分布和依赖结构分离建模（Sklar定理）。
        \item 可以独立选择边际分布和依赖结构（Copula函数），提供了比相关性更灵活的建模框架。
        \item 因此该陈述正确。
    \end{itemize}

    \textbf{D选项正确（陈述错误）。}
    \begin{itemize}
        \item 短期（3个月或更短）的相关性估计通常不稳定、波动较大。
        \item 原因：
        \begin{itemize}
            \item 样本量不足，估计误差大
            \item 相关性本身具有时变特征，短期内可能频繁变化
            \item 市场冲击和结构性变化会导致相关性突然跳跃
        \end{itemize}
        \item 因此，陈述"短期相关性估计通常非常稳定"是错误的。
    \end{itemize}

\end{explanation}

\checklist{
    掌握依赖结构特征（相关性只对特定分布有效，零相关不意味着独立，Copula可分别建模边际和依赖结构，短期相关性估计不稳定）
}


\begin{question}
    On December 31, 2008, Portfolio A had a market value of \$4,500,000. The historical standard deviation of daily returns was 1.8\%. Assuming that Portfolio A is normally distributed, calculate the daily VAR(2.5\%) on a dollar basis and state its interpretation. Daily VAR(2.5\%) is equal to:
\end{question}

\begin{choices}
    \item \$132,300, implying that daily portfolio losses will fall short of this amount 2.5\% of the time.
    \item \$111,375, implying that daily portfolio losses will only exceed this amount 2.5\% of the time.
    \item \$132,300, implying that daily portfolio losses will only exceed this amount 2.5\% of the time.
    \item \$111,375, implying that daily portfolio losses will fall short of this amount 2.5\% of the time.
\end{choices}

\begin{correctanswer}
    C
\end{correctanswer}

\begin{explanation}
    \textbf{步骤1：确定参数}
        \begin{itemize}
            \item 组合市值：\$4{,}500{,}000
            \item 日收益率标准差：$\sigma = 1.8\% = 0.018$
            \item 置信水平：2.5\%（尾部概率）
            \item 对应$z$值：$z_{0.025} = 1.96$（标准正态分布）
        \end{itemize}

    \textbf{步骤2：计算百分比VaR}
        \[
        \text{VaR}_{\%}(2.5\%) = z_{0.025} \times \sigma = 1.96 \times 0.018 = 0.03528 = 3.528\%
        \]

    \textbf{步骤3：计算美元VaR}
        \[
        \text{VaR}_{\$}(2.5\%) = \text{VaR}_{\%} \times \text{Portfolio Value} = 0.03528 \times \$4{,}500{,}000 = \$132{,}300
        \]

    \textbf{解释}
        \begin{itemize}
            \item VaR(2.5\%)表示：在2.5\%的概率下，每日组合损失会超过\$132{,}300
            \item 等价地：在97.5\%的概率下，每日组合损失不会超过\$132{,}300
        \end{itemize}
    

\end{explanation}

\checklist{
    掌握VaR计算和解释
}



\begin{question}
    Michael Wang is a risk analyst who wants to model the dependence between asset returns using copulas and must convince his manager that this is the best approach. Which of the following statements are correct?

    I. The dependence between the return distributions of portfolio assets is critical for risk measurement.

    II. Correlation estimates often appear stable in periods of low market volatility and then become volatile in stressed market conditions. Risk measures calculated using correlations estimated over long horizons will therefore underestimate risk in stressed periods.

    III. Pearson correlation is a linear measure of dependence between the return of two assets equal to the ratio of the covariance of the asset returns to the product of their volatilities.

    IV. Using copulas, one can construct joint return distribution functions from marginal distribution functions in a way that allows for more general types of dependence structure of the asset returns.
\end{question}

\begin{choices}
    \item I, II, and III
    \item II and IV
    \item I, II, III, and IV
    \item I, III, and IV
\end{choices}

\begin{correctanswer}
    D
\end{correctanswer}

\begin{explanation}
 \textbf{陈述I：正确。}
    \begin{itemize}
        \item 资产收益之间的依赖关系是风险度量的核心要素。依赖结构直接影响投资组合的分散化效果、尾部风险和整体风险水平。因此，准确建模依赖关系对风险度量至关重要。
    \end{itemize}

 \textbf{陈述II：错误。}
    \begin{itemize}
        \item 相关性在压力时期确实会上升且变得不稳定。但使用长期估计的相关性不一定总是"低估"风险。
        \item 影响取决于投资组合的具体结构（如做多/做空、相关性方向等）。在某些情况下，长期相关性可能反而会高估或产生其他偏差。因此，该陈述过于绝对化，不够准确。
    \end{itemize}

 \textbf{陈述III：正确。}
    \begin{itemize}
        \item Pearson相关系数确实是线性依赖度量。定义：$\rho_{X,Y} = \frac{\text{Cov}(X,Y)}{\sigma_X \sigma_Y}$，即协方差除以两个资产波动率的乘积。这是标准的线性相关性定义。
    \end{itemize}

 \textbf{陈述IV：正确。}
    \begin{itemize}
        \item Copula允许从边际分布函数构建联合分布函数。核心优势：可以独立选择边际分布和依赖结构（Sklar定理）。能够建模比线性相关性更一般的依赖类型（非线性、非对称、尾部依赖等）。
    \end{itemize}
\end{explanation}

\checklist{
    掌握Copula和相关性
}

\begin{question}
    Portfolios (X) and (Y) each have volatility of 30\%, but portfolio (Y) has a higher return and therefore its absolute VaR is lower; i.e., Absolute VaR = - return * T + deviate * volatility * $\sqrt{T}$. Which coherence property does this illustrate?
\end{question}

\begin{choices}
    \item Monotonicity
    \item Subadditivity
    \item Positive Homogeneity
    \item Translational invariance
\end{choices}

\begin{correctanswer}
    A
\end{correctanswer}

\begin{explanation}

    \textbf{A选项正确。}
    \begin{itemize}
        \item 单调性要求：如果投资组合A的收益总是优于B，则A的风险度量应小于等于B。
        \item 题目中：两个组合波动率相同（30\%），但组合Y收益更高。
        \item 根据绝对VaR公式：$\text{VaR} = -r \times T + z \times \sigma \times \sqrt{T}$，收益更高的组合Y具有更低的VaR。
        \item 这体现了单调性：更好的投资组合（更高收益）被度量出更低的风险。
    \end{itemize}

    \textbf{B选项错误。}
    \begin{itemize}
        \item 次可加性是关于组合合并后的风险应小于等于各自风险之和，与本题无关。
    \end{itemize}

    \textbf{C选项错误。}
    \begin{itemize}
        \item 正齐次性是关于风险随投资规模成比例变化，与本题无关。
    \end{itemize}

    \textbf{D选项错误。}
    \begin{itemize}
        \item 平移不变性是关于添加无风险资产时风险度量的调整，与本题无关。
    \end{itemize}

\end{explanation}

\checklist{
    掌握风险度量的一致性（单调性表示收益更高的组合风险更低，这是风险度量的基本性质）
}



\begin{question}
    An analyst at an investment bank is calculating the payoff of a correlation swap that is maturing. The notional amount of the correlation swap covers four stocks that are equally weighted. The analyst uses the realized pairwise correlations given below to calculate the correlation that will be applied to determine the payoff for the swap:
    \begin{table}[H]
        \centering
        \renewcommand{\arraystretch}{1.2}
        \begin{tabular}{|l|c|c|c|c|}
            \hline
            & \textbf{Stock A} & \textbf{Stock B} & \textbf{Stock C} & \textbf{Stock D} \\
            \hline
            \textbf{Stock A} & 1.0 & 0.5 & 0.7 & 0.6 \\
            \textbf{Stock B} & 0.5 & 1.0 & 0.3 & 0.4 \\
            \textbf{Stock C} & 0.7 & 0.3 & 1.0 & 0.8 \\
            \textbf{Stock D} & 0.6 & 0.4 & 0.8 & 1.0 \\
            \hline
        \end{tabular}
        \caption{股票相关系数矩阵}
    \end{table}
    What is the average realized correlation?
\end{question}

\begin{choices}
    \item 0.35
    \item 0.70
    \item 0.52
    \item 0.78
\end{choices}

\begin{correctanswer}
    B
\end{correctanswer}

\begin{explanation}


        \textbf{步骤1：确定参数}
        \begin{itemize}
            \item 股票数量：$n=4$
            \item 配对数量：$\binom{n}{2} = \frac{n(n-1)}{2} = \frac{4 \times 3}{2} = 6$
        \end{itemize}

        \textbf{步骤2：提取所有配对相关性（上三角或下三角，排除对角线）}
        \begin{itemize}
            \item $\rho(A,B) = 0.5$
            \item $\rho(A,C) = 0.7$
            \item $\rho(A,D) = 0.6$
            \item $\rho(B,C) = 0.3$
            \item $\rho(B,D) = 0.4$
            \item $\rho(C,D) = 0.8$
        \end{itemize}

    \textbf{步骤3：计算平均实现相关性}
        \[
        \bar{\rho} = \frac{2}{n(n-1)} \sum_{i<j} \rho_{i,j} = \frac{2}{4 \times 3} \times (0.5 + 0.7 + 0.6 + 0.3 + 0.4 + 0.8) = \frac{4.2}{6} = 0.70
        \]

        \textbf{说明}
        \begin{itemize}
            \item 相关性互换的平均相关性是所有配对相关性的简单平均。
            \item 公式：$\bar{\rho} = \frac{2}{n(n-1)} \sum_{i<j} \rho_{i,j}$，其中$i<j$表示所有不同的配对。
        \end{itemize}


\end{explanation}

\checklist{
    掌握相关性互换计算（平均实现相关性是所有配对相关性的算术平均值，用于计算互换支付）
}


\begin{question}
    Beta Capital Inc. (Beta), a wealth management company, currently manages an alternative investment fund, a growth equity fund, and individual assets that include Assets Z. The two funds are benchmarked against the same custom index. The alternative investment fund has allocated equal investments in the senior tranche and the mezzanine tranche of an MBS originated by Bank Alpha and issued by an SPV. Bank Alpha retains 6\% net economic interest in the securitization for the life of the transaction. To hedge its position, Beta purchased CDS on the MBS from Bank Gamma. Both Bank Alpha and Bank Gamma rely on retail deposits to finance their lending activities. Beta assesses its investment performance and capital budgeting using RAROC at the 99\% Confidence level. Beta's CIO expects market volatility and the correlation between assets in the custom index to increase over the coming year. The CIO then decides to execute the following two strategies:
    ·Enter into a payer volatility swap on Asset Z, where the prespecified fixed volatility is the current volatility of Asset Z.
    ·Take a long position in a correlation swap on the market index, with Beta as the fixed correlation payer.

    Assuming the CIO's expectations are realized, which of the following is correct?
\end{question}

\begin{choices}
    \item Assuming no low-risk anomaly, Beta would generate a gain from the volatility swap if the volatility of Asset Z decreases.
    \item If the realized correlation in the correlation swap is higher than the fixed correlation, Beta would generate a gain on the correlation swap.
    \item An increase in the default correlation between Bank Alpha and Bank Gamma, all else held constant, would lead to a decrease in the value of Beta's CDS position.
    \item Holding all else constant, Beta's estimated after-tax RAROC would not change if it adopts a 95-Confidence level compared to a 99-Confidence level for its capital budgeting.
\end{choices}

\begin{correctanswer}
    B
\end{correctanswer}

\begin{explanation}

    \textbf{A选项错误。}
    \begin{itemize}
        \item Beta是波动率互换的支付方（payer volatility swap），支付固定波动率。
        \item 支付函数：$\text{Payoff} = \text{Notional} \times (\sigma_{\text{realized}} - \sigma_{\text{fixed}})$
        \item 若实现波动率下降（$\sigma_{\text{realized}} < \sigma_{\text{fixed}}$），Beta需向对手方支付，产生亏损而非收益。
        \item Beta在实现波动率上升时获利，而非下降时。
    \end{itemize}

    \textbf{B选项正确。}
    \begin{itemize}
        \item Beta是相关性互换的固定相关性支付方（fixed correlation payer），同时持有多头头寸。
        \item 支付函数：$\text{Payoff} = \text{Notional} \times (\rho_{\text{realized}} - \rho_{\text{fixed}})$
        \item 若实现相关性 > 固定相关性，Beta作为固定支付方会收到净支付（多头获利）。
        \item CIO预期相关性增加，因此若实现相关性高于固定相关性，Beta将获利。
    \end{itemize}

    \textbf{C选项错误。}
    \begin{itemize}
        \item Beta购买了CDS对冲MBS头寸，CDS价值受MBS信用风险影响。
        \item 若Bank Alpha和Bank Gamma的违约相关性增加，意味着系统性风险上升，可能推高MBS的信用风险。
        \item 作为CDS买方，Beta的CDS保护价值上升（CDS价格上升），而非下降。
        \item 因此，违约相关性增加应导致CDS头寸价值增加，而非减少。
    \end{itemize}

    \textbf{D选项错误。}
    \begin{itemize}
        \item RAROC公式：$\text{RAROC} = \frac{\text{Expected Return}}{\text{Risk Capital}}$
        \item 置信水平从99\%降至95\%会降低风险资本要求（VaR/ES在更低置信水平下更小）。
        \item 即使预期收益不变，分母（风险资本）下降会导致RAROC上升。
        \item 因此，改变置信水平会影响RAROC，而非保持不变。
    \end{itemize}

\end{explanation}

\checklist{
    掌握互换交易策略（波动率互换支付方在波动率上升时获利，相关性互换多头在实现相关性高于固定相关性时获利）
}







\newpage




\section{考点2:相关性的经验特性}
\setcounter{question}{0}

\begin{question}
    A risk manager uses the past 500 months of correlation data from the S\&P 500 Index to estimate the long-run mean correlation of common stocks and the mean reversion rate. Based on historical data, the long-run mean correlation of S\&P 500 stocks was 35\%, and the regression output estimates the following regression relationship: Y = 0.225 – 0.80X. Suppose that in March 2023, the average monthly correlation for all S\&P 500 stocks was 40\%. What is the expected correlation for April 2023 assuming the mean reversion rate estimated in the regression analysis?
\end{question}

\begin{choices}
    \item 35\%
    \item 36\%
    \item 38\%
    \item 40\%
\end{choices}

\begin{correctanswer}
    B
\end{correctanswer}

\begin{explanation}

   
        \textbf{步骤1：确定参数}
        \begin{itemize}
            \item 长期均值相关性：$\mu = 35\%$
            \item 2023年3月相关性：$\rho_{t-1} = 40\%$
            \item 均值回归率：$\lambda = 0.80$（从回归系数$-0.80$得出）
        \end{itemize}

        \textbf{步骤2：计算2023年4月预期相关性}
        \begin{itemize}
            \item 均值回归模型：$\rho_t = \mu + \lambda(\rho_{t-1} - \mu)$
            \[
            \rho_t = 0.35 + 0.80 \times (0.40 - 0.35) = 0.35 + 0.04 = 0.36 = 36\%
            \]
        \end{itemize}

        \textbf{解释}
        \begin{itemize}
            \item 3月相关性（40\%）高于长期均值（35\%），偏离$+5\%$。
            \item 均值回归率80\%意味着偏离的80\%会在下一期回归：$0.80 \times 5\% = 4\%$。
            \item 因此，4月预期相关性 = $40\% - 4\% = 36\%$。
        \end{itemize}

\end{explanation}
\checklist{
        理解相关性均值回归（当前相关性会向长期均值回归）
}


\begin{question}
    Suppose a risk manager examines the correlations and correlation volatility of stocks in the
    S\&P 500 Index for the period beginning in 1980 and ending in 2020. Expansionary periods are defined as periods where the U.S. gross domestic product (GDP) growth rate is greater than 3.5\%, periods are normal when the GDP growth rates are between 0 and 3.5\%, and recessions are periods with two consecutive negative GDP growth rates. Which of the following statements characterizes correlation and correlation volatilities for this sample? The risk manager will most likely find that
\end{question}

\begin{choices}
    \item correlations and correlation volatility are highest for recessions.
    \item correlations and correlation volatility are highest for expansionary periods.
    \item correlations are highest for normal periods, and correlation volatility is highest for recessions.
    \item correlations are highest for recessions, and correlation volatility is highest for normal periods.
\end{choices}

\begin{correctanswer}
    D
\end{correctanswer}

\begin{explanation}
    \textbf{A选项错误。}
    \begin{itemize}
        \item 虽然衰退期相关性最高，但相关性波动性并非最高。
    \end{itemize}
    \textbf{B选项错误。}
    \begin{itemize}
        \item 扩张期相关性和波动性均最低。
    \end{itemize}
    \textbf{C选项错误。}
    \begin{itemize}
        \item 相关性最高出现在衰退期，而非正常期。
    \end{itemize}
    \textbf{D选项正确。}
    \begin{itemize}
        \item 1980年至2020年标普500指数股票月度相关性的研究发现，衰退期间相关性最高，正常时期相关性波动性最高。经济衰退和正常时期的相关波动率分别为82.5\%和85.2\%。
    \end{itemize}
\end{explanation}
\checklist{
    理解经济周期对相关性的影响（衰退期相关性高，正常期波动大）
}



\begin{question}
    Suppose mean reversion exists for a variable with a value of 35 at time period t-1.
    Assume that the long-run mean value for this variable is 45 and ignore the stochastic term
    included in most regressions of financial data. What is the expected change in value of the
    variable for the next period if the mean reversion rate is 0.5.
\end{question}

\begin{choices}
    \item -10
    \item -5
    \item 5
    \item 10
\end{choices}

\begin{correctanswer}
    C
\end{correctanswer}

\textbf{C选项正确。}
\begin{itemize}
    \item 均值回归率$\alpha$表示变量向长期均值回归的速度。
    \item 当前值$X_{t-1} = 35$，长期均值$\mu = 45$，偏离为$\mu - X_{t-1} = 45 - 35 = 10$。
    \item 下一期的预期变化：
    \[
    E[\Delta X_t] = \alpha \times (\mu - X_{t-1}) = 0.5 \times 10 = 5
    \]
    \item 因此，预期变化为$+5$，变量从$35$增加到$40$（向均值$45$回归）。
\end{itemize}
\checklist{
    理解均值回归（变量会以一定速率向长期均值回归）
}

\begin{question}
    In estimating correlation matrices, risk managers often assume an underlying distribution
    for the correlations. Which of the following statements most accurately describes the best-fit
    distributions for equity correlation distributions, bond correlation distributions, and default
    probability correlation distributions? The best-fit distribution for the equity, bond, and
    default probability correlation distributions, respectively are:
\end{question}

\begin{choices}
    \item lognormal, generalized extreme value, and normal.
    \item Johnson SB, generalized extreme value, and Johnson SB.
    \item beta, normal, and beta.
    \item Johnson SB, normal, and beta.
\end{choices}

\begin{correctanswer}
    B
\end{correctanswer}

\begin{explanation}
    \textbf{A选项错误。}
    \begin{itemize}
        \item 对数正态分布不适合描述股权相关性分布，正态分布不适合描述违约概率相关性分布。
    \end{itemize}
    \textbf{B选项正确。}
    \begin{itemize}
        \item 股权相关性分布和违约概率相关性分布最适合用Johnson SB分布描述，债券相关性分布最适合用广义极值分布描述。
    \end{itemize}
    \textbf{C选项错误。}
    \begin{itemize}
        \item Beta分布不适合描述股权相关性分布，正态分布不适合描述债券相关性分布。
    \end{itemize}
    \textbf{D选项错误。}
    \begin{itemize}
        \item 正态分布不适合描述债券相关性分布，Beta分布不适合描述违约概率相关性分布。
    \end{itemize}
\end{explanation}
\checklist{
    掌握相关性分布（股权和违约用Johnson SB，债券用极值分布）
}

\begin{question}
    A risk manager uses the past 500 months of correlation data from the S\&P 500 Index to estimate the long-run mean correlation of common stocks and the mean reversion rate. Based on historical data, the long-run mean correlation of S\&P 500 stocks was 35\%, and the regression output estimates the following regression relationship: Y = 0.225 – 0.80X. Suppose that in March 2023, the average monthly correlation for all S\&P 500 stocks was 40\%. What is the expected correlation for April 2023 assuming the mean reversion rate estimated in the regression analysis?
\end{question}

\begin{choices}
    \item 35\%
    \item 36\%
    \item 38\%
    \item 40\%
\end{choices}

\begin{correctanswer}
    B
\end{correctanswer}

\begin{explanation}

    \begin{explanation}

        \textbf{步骤1：确定参数}
            \begin{itemize}
                \item 长期均值相关性：$\mu = 35\%$
                \item 2023年3月相关性：$\rho_{t-1} = 40\%$
                \item 均值回归率：$\lambda = 0.80$（从回归系数$-0.80$得出）
            \end{itemize}
    
        \textbf{步骤2：计算偏离}
            \begin{itemize}
                \item 当前值偏离均值：$\rho_{t-1} - \mu = 40\% - 35\% = +5\%$
                \item 当前相关性高于长期均值，应向均值回归。
            \end{itemize}
    
        \textbf{步骤3：计算2023年4月预期相关性}
            \begin{itemize}
                \item 均值回归模型：$\rho_t = \mu + \lambda(\rho_{t-1} - \mu)$
                \[
                \rho_t = 0.35 + 0.80 \times (0.40 - 0.35) = 0.35 + 0.80 \times 0.05 = 0.35 + 0.04 = 0.36 = 36\%
                \]
            \end{itemize}
    
    
    \end{explanation}

\end{explanation}
\checklist{
        理解相关性均值回归（当前值会以一定速率向长期均值回归）
}


\newpage

\section{考点3:财务相关性建模}
\setcounter{question}{0}
\begin{question}
    Suppose a risk manager creates a copula function, C, defined by the equation:
    \begin{equation*}
        C\left[ G_1(u_1), \ldots, G_n(u_n) \right] = F_n \left[ F_1^{-1}(G_1(u_1)), \ldots, F_n^{-1}(G_n(u_n)); \rho_F \right]
        \end{equation*}

    Which of the following statements does not accurately describe this copula function?
\end{question}

\begin{choices}
    \item $G_i (u_i)$ are standard normal univariate distributions.
    \item $F_n$ is the joint cumulative distribution function.
    \item $F_1^{-1}$ is the inverse function of $F_n$ that is used in the mapping process.
    \item $\rho_F$ is the correlation matrix structure of the joint cumulative function $F_n$.
\end{choices}

\begin{correctanswer}
    A
\end{correctanswer}

\begin{explanation}
    \textbf{A选项正确。}
    \begin{itemize}
        \item $G_i (u_i)$是边际分布，不具有标准正态分布的特性。它们可以是任何形式的分布，不一定是标准正态分布。
    \end{itemize}
    \textbf{B选项错误。}
    \begin{itemize}
        \item $F_n$确实是联合累积分布函数，这是正确的描述。
    \end{itemize}
    \textbf{C选项错误。}
    \begin{itemize}
        \item $F_1^{-1}$确实是用于映射过程的逆函数，这是正确的描述。
    \end{itemize}
    \textbf{D选项错误。}
    \begin{itemize}
        \item $\rho_F$确实是联合累积分布函数$F_n$的相关矩阵结构，这是正确的描述。
    \end{itemize}
\end{explanation}
\checklist{
    理解Copula函数（边际分布不一定是标准正态分布）
}


\begin{question}
    In estimating correlation matrices, risk managers often assume an underlying distribution for the correlations. Which of the following statements most accurately describes the best-fit
    distributions for stock index correlation distributions, commodity correlation distributions, and credit spread correlation distributions? The best-fit distribution for the stock index, commodity, and credit spread correlation distributions, respectively are:
\end{question}

\begin{choices}
    \item lognormal, generalized extreme value, and normal.
    \item Johnson SB, generalized extreme value, and Johnson SB.
    \item beta, normal, and beta.
    \item Johnson SB, normal, and beta.
\end{choices}

\begin{correctanswer}
    B
\end{correctanswer}

\begin{explanation}
    \textbf{A选项错误。}
    \begin{itemize}
        \item 对数正态分布不适合描述股指相关性分布，正态分布不适合描述信用利差相关性分布。
    \end{itemize}
    \textbf{B选项正确。}
    \begin{itemize}
        \item 股指相关性分布和信用利差相关性分布最适合用Johnson SB分布描述，商品相关性分布最适合用广义极值分布描述。这是因为：
        \begin{itemize}
            \item Johnson SB分布能够很好地捕捉股指和信用利差相关性的非对称性
            \item 广义极值分布适合描述商品相关性的尾部风险
        \end{itemize}
    \end{itemize}
    \textbf{C选项错误。}
    \begin{itemize}
        \item Beta分布不适合描述股指相关性分布，正态分布不适合描述商品相关性分布。
    \end{itemize}
    \textbf{D选项错误。}
    \begin{itemize}
        \item 正态分布不适合描述商品相关性分布，Beta分布不适合描述信用利差相关性分布。
    \end{itemize}
\end{explanation}
\checklist{
    掌握相关性分布（股指和信用用Johnson SB，商品用极值分布）
}


\begin{question}
    Which of the following statements best describes a Student's t copula?
\end{question}

\begin{choices}
    \item A major disadvantage of a Student's t copula model is the transformation of the original
    marginal distributions in order to define the correlation matrix.
    \item The mapping of each variable to the new distribution is done by defining a mathematical
    relationship between marginal and unknown distributions.
    \item A Student's t copula maps the marginal distribution of each variable to the standard t
    distribution.
    \item A Student's t copula is seldom used in financial models because ordinal numbers are required.
\end{choices}

\begin{correctanswer}
    C
\end{correctanswer}

\begin{explanation}
    \textbf{A选项错误。}
    \begin{itemize}
        \item 转换边际分布不是Student's t copula的主要缺点，这是所有copula模型的共同特征。
    \end{itemize}
    \textbf{B选项错误。}
    \begin{itemize}
        \item 这种描述过于笼统，没有具体说明Student's t copula的映射特点。
    \end{itemize}
    \textbf{C选项正确。}
    \begin{itemize}
        \item Student's t copula通过将每个变量的边际分布按百分位数映射到标准t分布来构建。这种映射方式能够更好地捕捉金融数据的尾部相关性。
    \end{itemize}
    \textbf{D选项错误。}
    \begin{itemize}
        \item Student's t copula在金融模型中经常使用，特别是在需要捕捉尾部相关性的情况下。
    \end{itemize}
\end{explanation}
\checklist{
    理解Student's t copula（将边际分布映射到标准t分布）
}


\begin{question}
    Which of the following statements about correlation and copula are correct? 
    
    I. Copula allows us to model the dependence structure between variables independently of their
    marginal distributions. 
    
    II. Linear transformation of variables preserves their correlation
    structure.  
    
    III. Correlation can be effectively used to measure the relationship between variables with
    infinite variance. 
    
    IV. Correlation is an appropriate measure of dependence when the variables follow a multivariate normal distribution.
\end{question}

\begin{choices}
    \item I and IV only.
    \item II, III, and IV only.
    \item I and III only.
    \item II and IV only.
\end{choices}

\begin{correctanswer}
    A
\end{correctanswer}

\begin{explanation}
    \textbf{陈述I：正确。}
    \begin{itemize}
        \item Copula基于Sklar定理，可以将联合分布分解为边际分布和依赖结构。允许独立选择边际分布和依赖结构（Copula函数），提供比相关性更灵活的建模框架。
    \end{itemize}

    \textbf{陈述II：错误。}
    \begin{itemize}
        \item 线性变换（如$Y = aX + b$）对单一变量可能保持某些性质，但相关性结构可能改变。对不同变量应用不同的线性变换时，相关性通常会改变。只有特定的线性变换（如同时对所有变量应用相同的变换）才可能保持相关性不变。
    \end{itemize}

    \textbf{陈述III：错误。}
    \begin{itemize}
        \item 相关性的定义需要协方差和方差：$\rho = \frac{\text{Cov}(X,Y)}{\sigma_X \sigma_Y}$。
        \item 如果变量具有无限方差（如某些厚尾分布），方差$\sigma_X$或$\sigma_Y$不存在，相关性无法定义。
    \end{itemize}

    \textbf{陈述IV：正确。}
    \begin{itemize}
        \item 对于多元正态分布，相关性完全刻画了变量间的依赖关系。在正态分布假设下，相关性是衡量依赖关系的合适且充分的度量。
    \end{itemize}

\end{explanation}
\checklist{
    理解Copula和相关性（Copula可独立建模依赖结构，相关性要求有限方差）
}


\begin{question}
    A Gaussian copula is constructed to estimate the joint default probability of two assets
    within a one-year time period. Which of the following statements regarding this type of copula is incorrect?
    
\end{question}

\begin{choices}
    \item This copula requires that the respective cumulative default probabilities are mapped to a
    bivariate standard normal distribution.
    \item This copula defines the relationship between the variables using a default correlation
    matrix, ρM
    \item The term maps each individual cumulative default probability for asset i for time period ton
    a percentile-to-percentile basis.
    \item This copula is a common approach used in finance to estimate joint default probabilities
\end{choices}

\begin{correctanswer}
    B
\end{correctanswer}

\begin{explanation}

    \textbf{A选项错误（陈述正确）。}
    \begin{itemize}
        \item 高斯Copula确实需要将累积违约概率通过逆标准正态分布函数（分位数函数）映射到标准正态分布的变量。
        \item 然后使用双变量标准正态分布来建模联合违约概率。
    \end{itemize}

    \textbf{B选项正确（陈述错误）。}
    \begin{itemize}
        \item 对于两个资产，只需要一个相关系数$\rho$来描述它们之间的违约相关性，而非相关矩阵$\rho_M$。
        \item 相关矩阵通常用于多个资产（$n \geq 3$）的情况，需要$n \times n$的矩阵。
    \end{itemize}

    \textbf{C选项错误（陈述正确）。}
    \begin{itemize}
        \item 高斯Copula通过分位数-分位数映射（quantile-to-quantile mapping）将每个资产的累积违约概率转换为标准正态分布的分位数。具体过程：使用逆标准正态分布函数$\Phi^{-1}$将违约概率$p_i(t)$映射到标准正态变量$Z_i$。
    \end{itemize}

    \textbf{D选项错误（陈述正确）。}
    \begin{itemize}
        \item 高斯Copula是金融建模中估计联合违约概率的常用方法，特别是在信用风险建模和结构化产品定价中。
    \end{itemize}

\end{explanation}
\checklist{高斯Copula的参数特征}

\begin{question}
    A risk manager wants to study the behavior of a portfolio that depends on only two economic
    variables, X and Y. X is uniformly distributed between 5 and 9, and Y is uniformly distributed
    between 6 and 10. The risk manager wants to model their joint distribution, H(X,Y). The theorem
    of Sklar proves that, for any joint distribution H, there is a copula C such that:
\end{question}

\begin{choices}
    \item H(4X + 5, 4Y + 6) is equal to C[X,Y].
    \item H(X,Y) is equal to C[u,d] where u is the density marginal distribution of X and d is the
    density marginal distribution of Y.
    \item H(X,Y) is equal to C[(X – 5)/4, (Y – 6)/4].
    \item H[(X - 5)/4, (Y- 6)/4] is equal to C(X,Y).
\end{choices}

\begin{correctanswer}
    C
\end{correctanswer}

\begin{explanation}
    \textbf{A选项错误。}
    \begin{itemize}
        \item 变换后的变量范围超出了[0,1]区间，不符合copula的定义域要求。
    \end{itemize}
    \textbf{B选项错误。}
    \begin{itemize}
        \item copula使用的是累积分布函数，而不是密度函数。
    \end{itemize}
    \textbf{C选项正确。}
    \begin{itemize}
        \item 根据Sklar定理，对于连续边际分布$F_X(x) = u$和$F_Y(y) = v$，联合分布$H(x,y)$可以表示为：
        \[
        H(x,y) = C(F_X(x), F_Y(y)) = C(u, v)
        \]
        
        \item 对于均匀分布$U(a,b)$，累积分布函数为$F(x) = \frac{x-a}{b-a}$。
        
        \item 在这个例子中：
        \begin{itemize}
            \item $X \sim U(5, 9)$，因此$F_X(x) = \frac{x-5}{9-5} = \frac{x-5}{4}$，其中$u = \frac{x-5}{4} \in [0,1]$
            \item $Y \sim U(6, 10)$，因此$F_Y(y) = \frac{y-6}{10-6} = \frac{y-6}{4}$，其中$v = \frac{y-6}{4} \in [0,1]$
        \end{itemize}
        
        \item 因此，联合分布可以表示为：
        \[
        H(X, Y) = C\left(\frac{X-5}{4}, \frac{Y-6}{4}\right)
        \]
    \end{itemize}
    \textbf{D选项错误。}
    \begin{itemize}
        \item 变换顺序错误，应该是先计算累积分布函数，再应用copula。
    \end{itemize}
\end{explanation}
\checklist{
    理解Sklar定理（联合分布可通过copula和边际分布表示）
}




\end{document}